\documentclass[11pt,a4paper]{article}
\usepackage[UTF8]{ctex}
\usepackage{geometry}
\geometry{a4paper, margin=1in}
\usepackage{hyperref}
\usepackage{titlesec}
\usepackage{enumitem}
\usepackage{amsmath}

\title{CS6493 自然语言处理 \\ 项目进度报告 \\ \vspace{0.5em} \large 自动反思学术代理：面向学术文档分析的自评估对话系统}
\author{小组项目 - 课题三：使用 LlamaIndex 构建实用 LLM 应用}
\date{2026年3月}

\begin{document}

\maketitle

\section{引言与研究动机}

大型语言模型（LLM）的出现从根本上改变了人类与文本信息交互的方式。检索增强生成（Retrieval-Augmented Generation, RAG）代表了将 LLM 输出锚定于权威外部知识的最有前景的范式之一，从而缓解了广为人知的幻觉问题。然而，当应用于学术文献等专业领域时，传统 RAG 管道表现出显著的局限性：检索到的上下文可能相关性不足，生成的回答可能偏离源材料，且系统缺乏任何自我评估或质量控制机制。

本项目通过开发\textbf{自动反思学术代理}来解决这些局限性——这是一个专门为学术 PDF 文档设计的对话式问答系统。核心创新在于引入了自动化的自我评估和自我纠正机制。与直接向用户呈现生成回答的传统 RAG 系统不同，我们的代理包含一个内部"评论家"（Critic），在输出前评估回答质量。当初始回答未能达到预定义的质量标准时，系统会自主优化检索查询并重新生成改进的答案，形成一个无需人工干预即可运行的闭环反馈机制。

这种方法受到人工智能对齐研究最新进展的启发，特别是来自人工智能反馈的强化学习（RLAIF）和自反思 RAG（Self-RAG）框架，这些研究表明语言模型可以有效地评估和改进自身的输出。通过在实用的学术工具中实现这些原则，我们旨在弥合 LLM 自我改进的理论进展与学术研究工作流程中实际应用之间的差距。

\section{问题陈述与研究问题}

学术文档为问答系统带来了独特的挑战。研究论文包含密集的技术内容、专业术语、复杂的逻辑论证以及章节之间错综复杂的交叉引用。研究人员在阅读论文时可能提出从高层次摘要（"这项工作的主要贡献是什么？"）到具体技术细节（"所提方法如何处理边界情况？"）再到比较分析（"这种方法与之前的方法有何不同？"）等各类问题。

现有 RAG 系统在处理此类查询时往往表现不佳，原因有几个方面。首先，朴素的分块策略可能会割裂连贯的论述或分离相关概念，导致上下文检索不完整。其次，生成模型可能产生听起来合理但实际上错误的陈述——这在精确性至关重要的学术语境中是一种特别隐蔽的失败模式。第三，由于缺乏反馈机制，错误会静默传播，用户无法知道系统的回答是否可信。

我们的项目研究以下问题：
\begin{enumerate}
    \item 自动化的自我评估机制能否在没有人工监督的情况下有效识别低质量的 RAG 输出？
    \item 基于自生成反馈的迭代查询优化是否比单次生成能提高答案质量？
    \item 使用量化的本地 LLM 部署时，模型容量与回答质量之间存在怎样的权衡关系？
    \item 不同的文档分块策略如何影响下游的检索和生成性能？
\end{enumerate}

\section{方法论}

\subsection{系统架构概述}

所提出的系统采用模块化架构，包含四个相互连接的组件：文档处理、向量索引、回答生成和自我评估。工作流程如下：首先将学术 PDF 摄入并分割成语义连贯的文本块；然后将这些块嵌入到向量空间并建立索引以实现高效检索；当用户提出问题时，检索相关块并连同查询一起输入语言模型；最后，在呈现回答之前，内部评论家评估其质量，必要时触发优化流程。

\subsection{文档处理管道}

下游检索的质量关键取决于源文档的分割方式。我们实现并比较两种分块策略以研究它们的相对优劣。

第一种策略采用带重叠窗口的固定大小分块。文档被划分为预定令牌长度的片段（我们的实现中为 512 个令牌），相邻块共享小的重叠区域（50 个令牌）以保持边界处的局部上下文。这种方法计算效率高且确定性强，但可能会任意分割段落、公式或逻辑论证等连贯的文本单元。

第二种策略采用由嵌入相似度引导的语义分块。这种方法不施加严格的长度限制，而是识别语义内容显著变化的自然断点。通过计算连续句子的嵌入并检测相似度不连续性，系统可以保持自包含概念单元的完整性。虽然计算量更大，但假设这种策略能产生更好地服务于检索需求的高质量块。

\subsection{检索与对话上下文管理}

检索到的文本块构成回答生成的知识基础。我们采用针对学术文本微调的中英双语嵌入模型进行稠密向量检索。给定用户查询，系统从索引语料库中识别语义上最相似的前 $k$ 个块。

学术探究往往跨越多轮对话展开。用户可能首先询问论文的方法论，然后跟进关于具体实验细节的问题，使用预设共享上下文的代词和引用。为支持这种多轮交互，我们集成了对话记忆缓冲区来维护近期对话历史。这使系统能够解析指代并在延长的交互中保持主题连贯性。

\subsection{自反思生成循环}

我们系统的核心创新是自反思生成循环，它将传统的单次 RAG 管道转变为迭代自我改进过程。

收到用户查询后，系统首先执行标准 RAG：检索相关上下文并生成初始回答。然而，系统不是立即呈现这个回答，而是调用内部评论家模块。评论家沿着源自 RAG 三元组评估框架的两个维度评估草稿回答：

\textbf{忠实度}衡量生成的文本是否严格基于检索到的上下文。忠实的回答所做的声明可以直接追溯到源文本块中的特定段落，避免幻觉信息或无支持的概括。评论家根据明确的评估标准在 1-5 分范围内分配忠实度分数。

\textbf{答案相关性}衡量回答是否直接针对用户的查询。高度相关的回答提供所请求的具体信息，而没有过多的离题内容或回避性的模糊表述。这个维度同样在 1-5 分范围内评分。

评论家通过平均这两个维度来计算综合质量分数。如果此分数超过预定义的阈值（在我们的实现中设为 4.0），则回答被视为可接受并呈现给用户。否则，评论家生成诊断反馈，识别具体缺陷并提出可能产生更好上下文的优化查询。

然后系统使用这个优化后的查询重新检索上下文并生成新的回答，再进行另一轮评估。这个循环持续进行，直到达到质量阈值或达到最大迭代限制（设为 3 次），届时返回可用的最佳回答，并附带透明度指标显示发生了多少次优化循环。

\subsection{本地 LLM 部署}

为确保数据隐私并消除对外部 API 服务的依赖，我们使用 Ollama 框架在本地部署所有语言模型。Ollama 为量化模型提供高效推理，使得有能力的 LLM 能够在消费级硬件上部署。我们使用指令微调模型的 4 位量化版本，具体比较较大的模型（Qwen-2.5-7B-Instruct）与更紧凑的替代方案（Qwen-2.5-1.5B-Instruct），以表征质量-效率权衡。

\section{评估框架}

我们的评估方法完全聚焦于评估系统输出的语言和信息质量的自然语言处理指标。我们刻意排除了内存消耗或推理延迟等以硬件为中心的指标，因为这些是实现细节而非系统核心能力的衡量标准。

主要评估框架是 RAG 三元组，它全面评估任何 RAG 管道中的三个关键环节：

\textbf{上下文相关性}通过衡量检索到的块中有多大比例包含与回答查询相关的信息来评估检索质量。高上下文相关性表示高效检索且噪声最小。

\textbf{忠实度}通过衡量生成回答中有多大比例的声明可归因于检索到的上下文来评估生成质量。高忠实度表示准确、有据的生成，没有幻觉。

\textbf{答案相关性}通过衡量最终回答对用户实际信息需求的满足程度来评估端到端效用。高答案相关性表示系统成功理解了查询并提供了适当的信息。

除了标准的 RAG 三元组，我们引入了一个自定义指标来量化我们的核心创新：\textbf{自纠正成功率}。该指标衡量以下查询的比例：(a) 初始回答分数低于质量阈值，(b) 自反思循环被触发，且 (c) 优化后的最终回答超过了阈值。高自纠正成功率表明自动反馈机制在没有人工干预的情况下有效地提高了回答质量。

\section{当前进展}

开发按照三阶段计划推进：架构设计、实现和实验评估。截至本报告，前两个阶段已基本完成。

\subsection{已完成工作}

基础设施已全面投入运行。我们建立了包含所有必要依赖的 Python 开发环境，配置了用于本地模型部署的 Ollama 框架，并验证了量化的 Qwen 模型对我们的用例表现充分。

文档处理管道已实现，支持固定大小和语义两种分块策略。系统可以摄入 PDF 文件，提取文本内容，根据所选策略分割文档，并构建向量索引以实现高效检索。

核心自反思机制——我们的主要创新——已实现并通过单元测试验证。评论家模块成功解析结构化评估输出，计算综合质量分数，生成诊断反馈，并在初始回答不充分时提出优化查询。

使用 Streamlit 框架开发了交互式演示界面。该界面允许用户上传 PDF 文档，配置系统参数（模型选择、分块策略、质量阈值），提出问题，并检查包括中间评估和优化步骤在内的内部反思过程。

\subsection{剩余工作}

实验评估阶段将占据项目时间表的剩余部分。具体任务包括：

首先，基于选定的 NLP 研究论文构建包含 20-30 个不同复杂度问题（事实回忆、方法论解释、比较分析、多轮对话）的综合测试数据集。

其次，在各实验条件下执行系统性实验：两种 LLM 后端（7B 对比 1.5B）、两种分块策略（固定对比语义）、以及启用/禁用自反思机制。

第三，使用基于 LLM 的评估计算定量指标，并在样本子集上辅以人工验证。

第四，将研究发现综合到最终项目报告中，并准备课堂演示材料。

\section{结论}

本进度报告介绍了我们在自动反思学术代理方面的进行中工作，这是一个将自动自我评估和自我纠正引入 RAG 范式的对话系统。核心实现已完成：系统可以摄入学术 PDF，回答用户问题，并在初始输出低于质量标准时自主优化其回答。剩余工作聚焦于严格的实验评估，以量化这种方法的收益并表征其在不同配置下的行为。我们预期最终结果将证明自反思机制在提高专业领域 RAG 系统可靠性方面的实用价值。

\end{document}
